\section{Deterministic revealing optimization: the exact curve}
\label{sec:optional}

The goal of this section is to prove \cref{thm:optional-curve}, the rigorous
deterministic part of the informal revealing-optimization curve theorem
\cref{thm:intro-common-upper-curves}.  We now allow raw execution for a
finite common time $u$.  All jobs have this same known upper limit, and a
test reveals $p_j\in[0,u]$.  Thus the offline effective length of job $j$ is
\begin{equation}
 q_j=\min\{u,1+p_j\},
 \label{eq:opt:effective-length}
\end{equation}
and the offline optimum schedules the $q_j$ in nondecreasing order.  We write
$V(u)$ for the optimal deterministic size-asymptotic ratio at this fixed
common upper limit.

The answer has six branches.  We first define their transition points.  Put
\begin{equation}
 \varphi=\frac{1+\sqrt5}{2},
 \qquad
 u_\diamond=\frac{1+\sqrt{3+2\sqrt5}}2
 =1.866760399173862\ldots .
 \label{eq:opt:phi-udiamond}
\end{equation}
Let $s_0>1$ be the root of
\begin{equation}
 s_0^3=(s_0+1)^2,
 \qquad
 u_0=1+s_0=3.147899035704787\ldots .
 \label{eq:opt:u0}
\end{equation}
For $s=u-1$, put
\begin{equation}
 T_s(z)=sz^2+(s^3+s^2+1)z-(s^2+s+1)(s+2),
 \label{eq:opt:Ts}
\end{equation}
and let $\rho_I(u)$ be the positive root of
\begin{equation}
s\rho^2+(s^3+s^2+1)\rho-(s^2+s+1)(s+2)=0.
\label{eq:opt:rhoI-equation}
\end{equation}
Denote the left-hand side by $T_s(\rho)$.
Equivalently,
\begin{equation}
 \rho_I(u)=
 \frac{-1-u+2u^2-u^3+
 \sqrt{u^6-4u^5+10u^4-6u^3-3u^2+6u-3}}
 {2(u-1)}.
 \label{eq:opt:rhoI-explicit}
\end{equation}

At the other end, let $z_*>1$ solve $z_*-3=\log z_*$ and define
\begin{equation}
 r=\frac2{z_*-1},
 \qquad
 R_{\rm obl}=1+r=1.570574096649395\ldots,
 \qquad
 z_*=4.505241495792883\ldots .
 \label{eq:opt:obligatory-constant}
\end{equation}
For $\varphi+2\le u\le z_*$, define $C_{\rm mix}(u)=1+c$ by the unique
solution $(c,m)$ of
\begin{equation}
 \operatorname{artanh}m-m
 =1-\frac1c+\frac12\log\!\left(1+\frac2c\right),
 \qquad
 u=1+\frac2c-m,
 \label{eq:opt:mixed-parametrization}
\end{equation}
with $r\le c\le1/\varphi$ and $0\le m\le1/\varphi$.

\begin{theorem}[Exact revealing-optimization curve]
\label{thm:optional-curve}
For every fixed $u>0$, the optimal deterministic size-asymptotic ratio in
the revealing-optimization model is
\begin{equation}
\boxed{
V(u)=
\begin{cases}
1, &0<u\le1,\\[1mm]
u, &1\le u\le u_\diamond,\\[1mm]
\rho_I(u), &u_\diamond\le u\le u_0,\\[1mm]
\displaystyle 1+\frac1{\sqrt{u-1}},
   &u_0\le u\le\varphi+2,\\[3mm]
C_{\rm mix}(u), &\varphi+2\le u\le z_*,\\[1mm]
R_{\rm obl}, &u\ge z_*.
\end{cases}}
\label{eq:opt:full-curve}
\end{equation}
For every branch there is an explicit deterministic algorithm satisfying
\[
 \cost(\mathcal A,I)\le V(u)\cost(\OPT,I)+O_u(n).
\]
On finite $u\ge z_*$ the remainder is one absolute $O(n)$ term, and
\cref{thm:det-obligatory} supplies the same algorithm at $u=\infty$.
Conversely, for every deterministic algorithm $\mathcal A$ and every
$\epsilon>0$, there are arbitrarily large fixed instances $I$ such that
\[
 \cost(\mathcal A,I)
 \ge\bigl(V(u)-\epsilon\bigr)\cost(\OPT,I).
\]
\end{theorem}

The proof combines two kinds of lower bounds.  Binary hidden-stopping
instances give the first four branches, while a continuous capped version
of the harmonic construction gives the mixed branch and the plateau.  Three
explicit policies match them: raw execution, a forced-prefix uniform
threshold, and the history-dependent threshold from obligatory testing.
The upper-bound witness can be read from the following policy map.  The
algorithms themselves are specified formally in
\cref{alg:raw,alg:ute,alg:adaptive-threshold}.
\begin{center}
\begin{tabular}{@{}lll@{}}
\toprule
range & policy & parameter choice \\
\midrule
$0<u\le u_\diamond$
 & \textsc{Raw} & none \\
$u_\diamond<u\le u_0$
 & \textsc{ForcedPrefixUTE} & $b=b(u)$ from \eqref{eq:opt:ute-parameters} \\
$u_0<u\le\varphi+2$
 & \textsc{ForcedPrefixUTE} & $b=0$ \\
$\varphi+2<u<z_*$
 & \textsc{AdaptiveThreshold} & $c=c(u)$ from \eqref{eq:opt:mixed-parametrization} \\
$z_*\le u\le\infty$
 & \textsc{AdaptiveThreshold} & $c=r$ \\
\bottomrule
\end{tabular}
\end{center}

\subsection{The shape of the claimed curve}

The global-maximum assertion follows directly from the displayed formulas.
Only the algebraic branch needs a check.  Implicit differentiation of
$T_s(\rho_I(u))=0$, with $s=u-1$, gives
\begin{equation}
 \rho_I'(u)=
 -\frac{\rho^2+(3s^2+2s)\rho-3(s+1)^2}
 {2s\rho+s^3+s^2+1}.
 \label{eq:opt:rhoI-derivative}
\end{equation}
Using $T_s(\rho)=0$, its numerator factors as
\[
 \frac{2s^3+s^2-1}{s}
 \left(\rho-
 \frac{2s^3+3s^2-2}{2s^3+s^2-1}\right).
\]
On the relevant interval $s\ge u_\diamond-1>2/3$, the prefactor is
positive, while the fraction in parentheses is at most $5/3$ because
$4s^3-4s^2+1>0$.  The sign certificate
\eqref{eq:opt:T-five-thirds} gives $\rho>5/3$.  Hence
$\rho_I'(u)<0$.  The reciprocal branch and the mixed parametrization are
also decreasing, so $V$ attains its global maximum only at
$u=u_\diamond$.

The lower bounds, including the fact that the middle construction is safe
against raw execution, are proved next.  We then give self-contained upper
bounds through $u=\varphi+2$.  The parameterized upper bound for the last two
branches is proved separately in \cref{sec:optional-upper}.

\subsection{The binary stopping game and its lower envelope}
\label{sec:optional-lower}

The first four branches already occur on binary instances
$p_j\in\{0,u\}$.  The stopping point must be hidden from the algorithm;
announcing it in advance gives a strictly weaker lower bound.

\begin{lemma}[Raw-safe hidden stopping]
\label{lem:opt:hidden-stopping}
Fix $u>1$, a target $\rho\le u$, and a stopping fraction $\alpha\in(0,1)$.
Against a deterministic algorithm, there is a legal binary adversary whose
leading excess at its first crossing is
\begin{equation}
 (u-\rho)(1-\sigma^2)+\sigma^2 f_{u,\rho}(y,b),
 \label{eq:opt:stopping-decomposition}
\end{equation}
where $0\le b\le y\le1$ and
\begin{equation}
 0\le y-\alpha<\frac1{n-v}.
 \label{eq:opt:stopping-overshoot}
\end{equation}
Moreover,
\begin{equation}
 f_{u,\rho}(y,b)=
 1-\rho+2y+(u-1)(1-\rho)y^2
 +b^2+2b(u-1-uy).
 \label{eq:opt:stopping-f}
\end{equation}
Here $\sigma\in[0,1]$ is the fraction of jobs not already executed raw.
Consequently, if
\[
 \min_{0\le b\le\alpha}f_{u,\rho}(\alpha,b)\ge0,
\]
then the adversary gives ratio at least $\rho-O_u(1/n)$.
\end{lemma}

\begin{proof}
Start in a long phase.  A fresh job that is tested receives $p=u$, while a
fresh job run raw is assigned the hidden value zero.  If $v$ jobs have been
run raw and $L$ tested-long jobs have been seen, stop at the first touch for
which
\begin{equation}
 L\ge\alpha(n-v).
 \label{eq:opt:stopping-line}
\end{equation}
All untouched jobs are then assigned value zero.  Only for the purpose of
lower-bounding the remaining cost, we grant the algorithm this future
information and allow an optimal continuation; in the actual interaction,
values are still revealed only by tests.  If the line is never crossed, every job
was run raw: if even one job had been tested, then after the last first touch
the tested fraction would be one.  On the all-zero hidden instance the ratio
is $u\ge\rho$.

At a crossing, let $e$ of the $L$ long jobs already be complete, let $d=L-e$
be pending, and let $x=n-v-e-d$ jobs remain untouched.  The following
favorable canonical schedule lower-bounds the algorithm: run the $v$ raw
jobs; test and immediately process the $e$ completed long jobs; perform the
$d$ noncompleting long tests; test the $x$ zero jobs; and finally process the
$d$ pending long jobs.  Indeed, if a noncompleting unit test is immediately
before an available length-$u$ completion while $h$ jobs remain, the two
orders contribute $h+hu$ and $hu+(h-1)$, respectively; moving the
completion left saves one.  Raw completions are available from time zero,
and identical long jobs may be relabeled so that the $e$ completed ones are
tested first.  After the switch, a zero test-completion has length one and
therefore precedes a pending operation of length $u$.  These exchanges
respect precedence.

The exact costs of this relaxed schedule and of OPT are
\begin{align}
 A_{\rm ex}={}&
 u\left(vn-\frac{v(v-1)}2\right)
 +(1+u)\left(e(n-v)-\frac{e(e-1)}2\right) \notag\\
 &+d(n-v-e)+x(d+x)-\frac{x(x-1)}2
 +u\frac{d(d+1)}2,
 \label{eq:opt:stopping-Aexact}\\
 O_{\rm ex}={}&\frac{n(n+1)}2
 +(u-1)\frac{(e+d)(e+d+1)}2.
 \label{eq:opt:stopping-Oexact}
\end{align}
Put $\nu=v/n$, $\delta=(v+e+d)/n$, and $\lambda=e/n$.  Twice the quadratic
coefficients in \cref{eq:opt:stopping-Aexact,eq:opt:stopping-Oexact} are
\begin{align*}
 \mathcal A={}&1+2\delta(1-u\nu)+(u-1)\delta^2
 +2\nu(\nu+u-2)\\
 &+\lambda^2+2\lambda(\nu+u-1-u\delta),\\
 \mathcal O={}&1+(u-1)(\delta-\nu)^2.
\end{align*}
Now put $\sigma=1-\nu$, $y=(e+d)/(n-v)$, and $b=e/(n-v)$.
Direct expansion gives exactly \cref{eq:opt:stopping-decomposition} after
normalizing by $n^2/2$.  The first-crossing overshoot in
\cref{eq:opt:stopping-line} is less than one: a test increases
$e+d-\alpha(n-v)$ by one and a raw first touch increases it by $\alpha$.
This proves \cref{eq:opt:stopping-overshoot}, and in particular
$\sigma(y-\alpha)<1/n$.  There is one small domain issue: the actual
$b\le y$ can exceed $\alpha$.  Put $\bar b=\min\{b,\alpha\}$.  Then
$0\le b-\bar b\le y-\alpha$.  Since $f_{u,\rho}$ is Lipschitz on the unit
square,
\begin{equation}
 \sigma^2 f_{u,\rho}(y,b)
 \ge \sigma^2 f_{u,\rho}(\alpha,\bar b)-O_u(1/n).
\end{equation}
Thus the stated minimum over $0\le\bar b\le\alpha$ suffices uniformly even
when $n-v$ is small.  The discarded diagonal terms are $O_u(n)$ in the
unnormalized costs.
Finally, the adaptive transcript defines a fixed instance by replay: every
tested answer was fixed when revealed, and every raw job's zero value stayed
hidden.  This proves the lemma.
\end{proof}

Minimizing \cref{eq:opt:stopping-f} over the algorithm's completion choice
gives
\begin{equation}
 b_*(y)=\max\{0,uy-(u-1)\}.
 \label{eq:opt:bstar}
\end{equation}
The upper constraint $b\le y$ never binds, because
$uy-(u-1)\le y$ for $y\le1$.
There are therefore two elementary quadratic branches.  On the positive
branch, put
\begin{equation}
 A_u=u^2-u+1,
 \qquad
 D_{u,\rho}=A_u+(u-1)\rho.
\end{equation}
Then
\begin{equation}
 \min_b f_{u,\rho}(y,b)
 =1-\rho-(u-1)^2+2A_uy-D_{u,\rho}y^2.
 \label{eq:opt:stopping-positive-branch}
\end{equation}
Its maximizer is $\alpha=A_u/D_{u,\rho}$, and tangency to zero is
\begin{equation}
 \bigl(1-\rho-(u-1)^2\bigr)D_{u,\rho}+A_u^2=0.
 \label{eq:opt:stopping-tangency}
\end{equation}
Expanding \cref{eq:opt:stopping-tangency} gives
\cref{eq:opt:rhoI-equation}.  On the zero-$b$ branch,
\begin{equation}
 \min_b f_{u,\rho}(y,b)
 =1-\rho+2y+(u-1)(1-\rho)y^2,
\end{equation}
and tangency gives
\begin{equation}
 \rho=1+\frac1{\sqrt{u-1}},
 \qquad
 \alpha=\frac1{\sqrt{u-1}}.
 \label{eq:opt:zero-b-tangency}
\end{equation}
The latter point belongs to the zero-$b$ branch iff
$(u-1)^3\ge u^2$, with equality at $u=u_0$.

For $\rho=u$, the maximum in
\cref{eq:opt:stopping-positive-branch} is nonnegative exactly while
\begin{equation}
 [u(u-1)]^2-u(u-1)-1\le0,
\end{equation}
whose positive equality is $u=u_\diamond$.  We have proved the following
binary lower envelope.

\begin{proposition}[Binary lower curve]
\label{prop:opt:binary-lower}
Every deterministic algorithm has size-asymptotic ratio at least
\begin{equation}
 B(u)=
 \begin{cases}
 1,&0<u\le1,\\
 u,&1\le u\le u_\diamond,\\
 \rho_I(u),&u_\diamond\le u\le u_0,\\[1mm]
 \displaystyle1+\frac1{\sqrt{u-1}},&u\ge u_0.
 \end{cases}
 \label{eq:opt:binary-curve}
\end{equation}
The lower bound uses only $p_j\in\{0,u\}$.
\end{proposition}

\begin{proof}
For $1<u\le u_\diamond$, use $\rho=u$ and the maximizing point of
\cref{eq:opt:stopping-positive-branch}.  For
$u_\diamond\le u\le u_0$, use the positive root $\rho_I(u)$ and
$\alpha=A_u/D_{u,\rho_I(u)}$; the tangency identity makes the maximum zero.
The minimizing completion mass is positive in the interior and becomes zero
at $u=u_0$.  For $u\ge u_0$, use the zero-$b$ tangency point from
\cref{eq:opt:zero-b-tangency}.  The branch conditions are precisely the two
inequalities checked above.  Finally, $V(u)\ge1$ is trivial for $u\le1$.
In each nontrivial case \cref{lem:opt:hidden-stopping} supplies fixed binary
instances by replay.
\end{proof}

\subsection{Matching upper bounds through the binary regime}
\label{sec:opt:low-upper}

The first policy is formalized in \cref{alg:raw}.

\begin{algorithm}[H]
\caption{\textsc{Raw}$(J,u)$}
\label{alg:raw}
\begin{algorithmic}
\Statex \textbf{Input:} jobs $J$ and a finite common upper limit $u$.
\State Fix an arbitrary order $j_1,\ldots,j_n$.
\State \textbf{For} $i=1,\ldots,n$, execute $j_i$ untested for $u$ time.
\end{algorithmic}
\end{algorithm}

For $u\le u_\diamond$, use \textsc{Raw}.  It is optimal when $u\le1$;
when $u>1$, every effective length in
\cref{eq:opt:effective-length} is at least one, so the ratio is at most $u$.

For the next interval, let $\rho=\rho_I(u)$, $s=u-1$, and define
\begin{equation}
 K=s^2+s+1,
 \quad D=K+s\rho,
 \quad b=\frac{K-s^2\rho}{D},
 \quad \chi=\frac{1-b}{\rho},
 \quad \alpha=\frac KD.
 \label{eq:opt:ute-parameters}
\end{equation}
The defining equation for $\rho$ gives
\begin{equation}
 0\le b<\chi<\alpha<1,
 \qquad
 \chi=\frac{s(s+1)}D,
 \qquad
 b=(s+1)\alpha-s;
 \label{eq:opt:ute-order}
\end{equation}
For completeness, these inequalities follow without numerical root
isolation.  Since $T_s$ is increasing on positive arguments,
\begin{equation}
 T_s\!\left(\frac K{s^2}\right)
 =\frac K{s^3}\bigl(-s^3+s^2+2s+1\bigr)\ge0
 \qquad(s\le s_0),
\end{equation}
and hence $b=(K-s^2\rho)/D\ge0$, with equality only at $s=s_0$.  Moreover,
\begin{equation}
 \alpha-\chi=\frac1D,
 \qquad
 \chi-b=\frac{s^2\rho-1}{D}>0.
\end{equation}
For the last sign,
\begin{equation}
 9T_s(5/3)=f(s):=6s^3-12s^2-2s-3<0.
 \label{eq:opt:T-five-thirds}
\end{equation}
throughout the present interval.  Indeed, on $[4/5,1]$ the polynomial is
decreasing and $f(4/5)=-1151/125$; on $[1,s_0]$ it has one interior minimum,
while $f(1)=-11$ and $f(s_0)<f(13/6)=-95/36$.  Here
$s_0<13/6$ follows by substituting $13/6$ in
$s^3-(s+1)^2$.  Thus $\rho>5/3$, while $s>4/5$, and
$s^2\rho>16/15$.  Finally $K<D$ gives $\alpha<1$.  In particular, $b=0$
occurs only at $u=u_0$.
The remaining two elementary branches use the single procedure in
\cref{alg:ute}.

\begin{algorithm}[H]
\caption{\textsc{ForcedPrefixUTE}$(J,u,b)$}
\label{alg:ute}
\begin{algorithmic}
\Statex \textbf{Input:} jobs $J$, finite $u>1$, and $b\in[0,1]$.
\State Fix a test order $j_1,\ldots,j_n$; set
       $k\gets\lfloor bn\rfloor$, $\tau\gets\min\{1,u-1\}$,
       and $Q\gets\varnothing$.
\State \textbf{For} $i=1,\ldots,n$ \textbf{do}:
\State \quad Test $j_i$ for one unit and observe $p_{j_i}$.
\State \quad \textbf{If} $i\le k$ or $p_{j_i}\le\tau$, process $j_i$
       immediately.
\State \quad \textbf{Otherwise}, add $j_i$ to $Q$.
\State Process the jobs of $Q$ in nondecreasing revealed processing time.
\end{algorithmic}
\end{algorithm}

On $u_\diamond\le u\le u_0$, use
\textsc{ForcedPrefixUTE} with $b$ from \cref{eq:opt:ute-parameters}.  Rounding
$bn$ changes the cost by $O_u(n)$.

Here is the complete binary calculation.  Adjacent exchanges place as many
long jobs as possible in the forced prefix and put all suffix-long jobs
before suffix zeros.  If $x$ is the total long mass, then the prefix-long
mass is $a=\min\{b,x\}$ and the deferred-long mass is
$d=(x-b)_+$.  After division by $n^2$, the leading excess
$G(x)=A-\rho O$ is
\begin{equation}
G(x)=
\begin{cases}
\displaystyle
 \frac{1-\rho}{2}+u\left(x-\frac{x^2}{2}\right)
 -\frac{\rho s}{2}x^2,
 &0\le x\le b,\\[3mm]
\displaystyle
 \frac{1-\rho}{2}
 +u\left(b-\frac{b^2}{2}\right)
 +(1-b)(x-b)+\frac{s}{2}(x-b)^2
 -\frac{\rho s}{2}x^2,
 &b\le x\le1.
\end{cases}
\label{eq:opt:binary-upper-gap}
\end{equation}
The second quadratic is strictly concave, has its stationary point at
$x=\alpha$, and satisfies $G(\alpha)=0$; after substituting
\cref{eq:opt:ute-parameters}, the latter identity is exactly
\cref{eq:opt:rhoI-equation}.  Since $b<\alpha$, its right derivative at $b$
is positive.  The left derivative satisfies
$G'_-(b)=G'_+(b)+s(1-b)>0$, and the first quadratic is concave, so it is
increasing on $[0,b]$.  Thus $G(x)\le0$ for every $x\in[0,1]$, including
the limiting case $b=0$.  This proves the binary upper bound without a
hidden local-versus-global step.

For $u_\diamond\le u\le2$, arbitrary values reduce to binary values by
independent Bernoulli extremalization, in the following order.  First raise
\emph{every} suffix-deferred value to $u$.  This cannot lower ALG, while its
offline effective length was already $u$.  Only after this deterministic
step, replace a prefix value $p$ by $u$ with probability $p/u$ and by zero
otherwise, and replace a suffix-immediate value $p\le u-1$ by $u$ with
probability $p$ and by zero otherwise.  Notice that the $u$ realization
of a suffix-immediate job becomes deferred, so status is not preserved.
Nevertheless the online check is pairwise.  For jobs $i$ before $j$, the
leading pair contribution is
\begin{equation}
 g_{ij}=\begin{cases}
 1+p_i,&i\text{ is immediate},\\
 2+p_j,&i\text{ is deferred and }j\text{ is immediate},\\
 2+\min\{p_i,p_j\},&i,j\text{ are deferred}.
 \end{cases}
 \label{eq:opt:pair-contribution}
\end{equation}
A rounded prefix job stays immediate and preserves its expected value.  If
$i$ is suffix-immediate, its rounded contribution is one when it becomes
zero and at least two when it becomes $u$, hence its expectation is at least
$1+p_i$.  If $i$ is deferred and $j$ is suffix-immediate, then $i$ has been
raised to $u$ and the expected pair contribution is
$2+up_j\ge2+p_j$; a deferred--deferred pair can only increase.  Thus
expected ALG does not decrease, up to its $O_u(n)$ diagonal terms.

The expected effective lengths are,
respectively,
\begin{equation}
 1+\frac{u-1}{u}p\le q(p),
 \qquad
 1+(u-1)p\le1+p,
\end{equation}
and $u$ for a deferred value.  Since the minimum is jointly concave,
$\mathbb E\min(q_i',q_j')\le\min(q_i,q_j)$.  Thus some binary realization
has competitive excess at least that of the original instance.

It remains to justify arbitrary values for $2\le u\le u_0$.  The following
quadratic lemma records the entire endpoint calculation.

\begin{lemma}[UTE endpoint game]
\label{lem:opt:ute-endpoint-game}
For $2\le u\le u_0$, UTE with the parameters
\cref{eq:opt:ute-parameters} satisfies
\begin{equation}
 \ALG\le\rho_I(u)\OPT+O_u(n).
\end{equation}
\end{lemma}

\begin{proof}
Fix the symbolic prefix/immediate/deferred decisions.  Convexity in one
coordinate reduces a prefix job to $p\in\{0,u\}$, a suffix-immediate job to
$p\in\{0,1\}$, and a suffix-deferred job to $p\in\{1^+,u\}$.  For the last
claim, on $1<p<u-1$ the deferred--deferred term has coefficient
$1-\rho<0$ multiplying a concave minimum; on $[u-1,u]$, OPT is flat and ALG
is nondecreasing.  Adjacent exchanges give the canonical endpoint order.

Normalize counts by $n$.  Let $a$ be prefix-$u$ mass, $d$ suffix-$u$ mass,
$t$ suffix-$1^+$ mass, and $m$ suffix-immediate-one mass.  Their constraints
are
\begin{equation}
 0\le a\le b,
 \qquad d,t,m\ge0,
 \qquad d+t+m\le1-b.
\end{equation}
Pairwise accounting using \cref{eq:opt:pair-contribution} gives
\begin{align}
 A={}&\frac12+(s+1)\left(a-\frac{a^2}{2}\right)
 +(1-b)(d+t+m)-\frac{m^2}{2}+\frac{s}{2}d^2,
 \label{eq:opt:ute-A}\\
 O={}&\frac12+\frac{(a+d+t+m)^2}{2}
 +\frac{s-1}{2}(a+d)^2.
 \label{eq:opt:ute-O}
\end{align}
Moving $m$ to $t$ does not change $O$ and increases $A$, so $m=0$.  Put
$x=a+d$ and write $F(x)=A-\rho O$ after optimizing the other variables.
The derivative of the gap in $t$ is
$1-b-\rho(x+t)$, so its maximizing value is $(\chi-x)_+$.  At fixed
$x,t$, moving mass from $d$ to $a$ leaves $O$ unchanged and changes $A$ at
rate $s(1-x)+b-a\ge0$.  Hence
\begin{equation}
 t=(\chi-x)_+,
 \qquad
 a=\min\{b,x\},
 \qquad d=(x-b)_+.
 \label{eq:opt:ute-reduction}
\end{equation}

For $x\ge\chi$, the gap $A-\rho O$ is the binary concave quadratic already
discussed; it is tangent to zero at $x=\alpha$.  For $x<\chi$, put
\begin{equation}
 C_0=-\frac{\rho-1}{2}+\frac{(1-b)^2}{2\rho},
 \qquad
 \gamma=s-\rho(s-1).
\end{equation}
Substitution gives
\begin{equation}
 A-\rho O
 =C_0+a(s+b-sx)-\frac{a^2}{2}+\frac{\gamma x^2}{2}.
 \label{eq:opt:ute-positive-t}
\end{equation}
One has $\gamma>0$ on $1\le s\le s_0$.  Indeed,
$T_s(2)=s(s^2-s+1)>0$, so for $s\le2$ the equation $T_s(\rho)=0$ gives
$\rho<2$, and hence $\gamma>2-s\ge0$.  For
$2\le s\le s_0<13/6$,
\begin{equation}
 16T_s(7/4)=12s^3-20s^2+s-4>0;
\end{equation}
the polynomial is $14$ at $s=2$ and has positive derivative thereafter.
Thus $\rho<7/4$ and $\gamma>7/4-3s/4>1/8$.

On $0\le x\le b$, one has $a=x$, and
\begin{equation}
 A-\rho O=C_0+(s+b)x
 -\frac{s+1+\rho(s-1)}2x^2.
 \label{eq:opt:ute-small-x}
\end{equation}
Its derivative is decreasing, and the derivative at $b$ multiplied by $D$
is
\begin{equation}
 \rho\bigl(1+s^2+\rho s^2(s-1)\bigr)>0,
\end{equation}
so the maximum is at $b$.  On $b\le x\le\chi$,
\cref{eq:opt:ute-positive-t} is convex and its maximum is at $b$ or $\chi$.
The endpoint $\chi$ lies on the binary branch.  Finally,
\begin{equation}
 F(b)-F(\chi)
 =\frac{\chi-b}{2}\,[2sb-\gamma(b+\chi)],
\end{equation}
and direct use of \cref{eq:opt:ute-parameters} gives
\begin{equation}
 D[\gamma(b+\chi)-2sb]
 =-s+(1+s-s^3)\rho+s^2(s-1)\rho^2>0.
 \label{eq:opt:ute-small-certificate}
\end{equation}
For a rational sign certificate, $T_s$ is increasing on positive arguments
and
\begin{equation}
 9T_s(5/3)=6s^3-12s^2-2s-3<0,
\end{equation}
so $\rho>5/3$.  If the right-hand side of
\cref{eq:opt:ute-small-certificate} is denoted by $q_s(\rho)$, then
\begin{equation}
 3\,\partial_\rho q_s(5/3)=7s^3-10s^2+3s+3>0.
\end{equation}
The last polynomial is increasing for $s\ge1$, so $q_s$ is increasing for
$\rho\ge5/3$.  Finally,
\begin{equation}
 9q_s(5/3)=10s^3-25s^2+6s+15>0.
\end{equation}
For the last sign, write $s=3/2+z$; the expression becomes
$3/2-3z/2+20z^2+10z^3$.  It is at least $3/2$ for
$-1/2\le z\le0$, and for $z\ge0$ it is larger than
$3/2-9/320$.  These inequalities prove that every endpoint gap is
nonpositive.  Diagonals and prefix rounding contribute $O_u(n)$.
\end{proof}

At $u=u_0$, $b=0$.  The same zero-prefix algorithm remains exact through
$u=\varphi+2$.

\begin{lemma}[Zero-prefix extension]
\label{lem:opt:zero-prefix}
For $u_0\le u\le\varphi+2$,
\textnormal{\textsc{ForcedPrefixUTE}}$(J,u,0)$ from \cref{alg:ute} is
$\bigl(1+1/\sqrt{u-1}\bigr)$-competitive up to $O_u(n)$.
\end{lemma}

\begin{proof}
The same coordinatewise reduction leaves capped-deferred mass $d$,
boundary-deferred mass $t$, immediate-one mass $m$, and zeros.  With
$s=u-1$, these masses satisfy $d,t,m\ge0$ and $d+t+m\le1$, and the leading
costs are
\begin{equation}
 A=\frac12+d+t+m-\frac{m^2}{2}+\frac{s d^2}{2},
 \qquad
 O=\frac12+\frac{(d+t+m)^2}{2}+\frac{(s-1)d^2}{2}.
\end{equation}
Move $m$ to $t$ and put $y=d+t$.  For fixed $y$, the ratio is
fractional-linear in $d^2$, so its maximum is on one of two faces:
\begin{equation}
 \frac{1+2y}{1+y^2}\le\varphi
 \quad(d=0),
 \qquad
 1+\frac{2y}{1+s y^2}\le1+\frac1{\sqrt s}
 \quad(d=y).
\end{equation}
The first one-variable function is maximized at $y=1/\varphi$ and the second
at $y=1/\sqrt s$, which gives the two displayed constants.
The second value dominates the first exactly when
$u\le\varphi+2$.  Binary instances attain it.
\end{proof}

\subsection{The unrestricted mixed lower bound}
\label{sec:opt:mixed-lower}

For $\varphi+2<u<z_*$, neither the binary family nor the pure cap-free
family is tight.  The correct adversary puts a capped block before a scaled
harmonic core.  Two lemmas ensure that optional raw
execution and early processing of caps cannot escape it.

The cap-free construction needed below is exactly the scaled harmonic core
from \cref{lem:opt:harmonic-block}; no second phase-accounting proof is
needed here.  We now prepend capped jobs to that reusable core.

\begin{lemma}[Cap deferral]
\label{lem:opt:cap-deferral}
Prepend capped mass $m$ with value $p=u$ to a later cap-free block of mass
$s=1-m$ and total processing mass $L$.  Suppose every later value satisfies
$1+p\le u$ and
\begin{equation}
 s(u-1)-L-m\ge0.
 \label{eq:opt:cap-deferral-condition}
\end{equation}
Then processing any capped job before all later jobs are complete cannot
decrease the leading online cost.
\end{lemma}

\begin{proof}
A cap completed after its test block but before a later job of value $p$
has mutual cost $1+u$; leaving it pending until after that job gives $2+p$.
The exchange saves $u-1-p\ge0$.  If cap mass $q$ is instead completed inside
the cap-test block, its pairs with all later jobs add
$q[s(u-1)-L]$.  Its largest possible saving against not-yet-tested caps is
$qm-q^2/2$.  Relative to leaving every cap pending, the change is therefore
at least
\begin{equation}
 q[s(u-1)-L-m]+\frac{q^2}{2}\ge0.
\end{equation}
The finite diagonal discrepancy is lower order.
\end{proof}

After cap deferral, pairwise accounting gives the offline coefficient $O$
and the guaranteed online lower benchmark $A^{\rm lb}$:
\begin{align}
 A^{\rm lb}={}&A_0^{\rm lb}+m(1-m)+m(1+L)+\frac u2m^2,
 \label{eq:opt:mixed-A}\\
 O={}&O_0+m(1-m+L)+\frac u2m^2.
 \label{eq:opt:mixed-O}
\end{align}
Indeed, a cap--later pair contributes $2+p$ online and $1+p$ offline,
whereas a cap--cap pair contributes $2+u$ online and $u$ offline; collecting
these pairs yields the two cross terms and the cap squares above.
In particular, $A^{\rm lb}-O=(A_0^{\rm lb}-O_0)+m$.

\begin{lemma}[Raw-safe quota and replay]
\label{lem:opt:mixed-raw-safe}
Fix $m\in(0,1)$ and a later cap-free block satisfying the hypotheses of
\cref{lem:opt:cap-deferral}.  Against every deterministic optional algorithm
there are fixed instances whose asymptotic excess is at least the cap-free
hybrid benchmark $A^{\rm lb}-CO$ from
\cref{eq:opt:mixed-A,eq:opt:mixed-O} for every target
$C\le u$.
\end{lemma}

\begin{proof}
In the initial phase, a tested fresh job receives $p=u$, while a job run raw
is assigned hidden value zero.  If $v$ raw jobs and $\ell$ tested caps have
been touched, stop at the first crossing
\begin{equation}
 \frac{\ell}{n-v}\ge m.
 \label{eq:opt:mixed-quota}
\end{equation}
With $S=n-v$, the overshoot is less than one, so $\ell=mS+O(1)$.  If no
crossing occurs, no test can have occurred: after all first touches, any
positive $\ell$ would make the ratio in \cref{eq:opt:mixed-quota} equal one.
Thus every job was run raw and the all-zero hidden instance has ratio
$u\ge C$.

At a crossing, scale the prescribed later instance to the remaining
$S-\ell$ first touches, absorbing the one-job cap overshoot and integer
rounding into its zero group.  A later test reveals the next value in this
fixed type stream; a later raw action consumes the next value but keeps it
hidden from the optional algorithm.

Move every initial raw-zero completion to the start.  It has the same length
$u$ as cap processing, and crossing a unit test only lowers completion cost.
With
\begin{equation}
 Z_v=vn-\frac{v(v-1)}2,
\end{equation}
this prefix contributes $uZ_v$ online and $Z_v$ offline.  Delete it.  The
remaining $S$ labels consist, up to one-job rounding, of capped fraction $m$
and the prescribed later cap-free block.

If the optional algorithm later runs a fresh assigned value $p$ raw, a
physical cap-free simulation tests and immediately processes that job in
time $1+p\le u$, suppresses the answer from the virtual optional algorithm,
and idles until its virtual $u$-block ends.  Inductively every physical
completion is no later than its virtual optional completion.  Since the
finite-cap and cap-free offline lengths are both $1+p$ on the later block,
\begin{equation}
 \ALG_{\rm opt}-C\OPT_{\rm opt}
 \ge (u-C)Z_v+
 \bigl(\ALG_{\rm hyb}-C\OPT_{\rm hyb}\bigr)-O_u(n).
 \label{eq:opt:mixed-raw-decomposition}
\end{equation}
The first term is nonnegative.  Finally, record every assigned value on its
job label.  Replaying this fixed vector produces the same transcript because
the values learned only by the physical simulation were hidden from the
virtual algorithm.  This proves legality.
\end{proof}

Apply \cref{lem:opt:harmonic-block} with $s=1-m$.  Substitution into
\cref{eq:opt:mixed-A,eq:opt:mixed-O} gives the guaranteed two-parameter
lower ratio
\begin{equation}
 \mathcal R_u(m,t)=1+
 \frac{m+\frac{(1-m)^2}{2}(1-t^{-2})}
 {\frac{(1-m)^2}{4}(1+t^{-2}+2\log t)
  +m(1-m)(1+\log t)+\frac u2m^2}.
 \label{eq:opt:mixed-ratio}
\end{equation}
At an interior stationary point with value $1+c$, differentiation and the
zero-gap equation give
\begin{equation}
 \log t=u-2-\frac1c,
 \qquad
 t^2=\frac{1-m}{1+m}\left(1+\frac2c\right).
 \label{eq:opt:mixed-stationarity}
\end{equation}
Eliminating $t$ gives precisely
\cref{eq:opt:mixed-parametrization}.  Conversely, those equations make
$\mathcal R_u(m,t)=1+c$.

The branch runs from
\begin{equation}
 (u,m,t)=\left(\varphi+2,\frac1\varphi,1\right)
 \quad\text{to}\quad
 (u,m,t)=(z_*,0,\sqrt{z_*}).
 \label{eq:opt:mixed-endpoints}
\end{equation}
It satisfies $u\ge\varphi+2$, $\log t<1$, and $m\le1/\varphi$.  Therefore
\begin{equation}
 (1-m)(u-1-\log t)
 >(1-m)\varphi\ge\frac1\varphi\ge m,
\end{equation}
so \cref{eq:opt:cap-deferral-condition} holds; also $1+p<3<u$ for every
later positive value in a sufficiently fine finite construction.  The
functions
\begin{equation}
 g(c)=1-\frac1c+\frac12\log\!\left(1+\frac2c\right),
 \qquad h(m)=\operatorname{artanh}m-m
\end{equation}
have $g'(c)=2/[c^2(c+2)]>0$ and
$h'(m)=m^2/(1-m^2)>0$.  Hence $m(c)=h^{-1}(g(c))$ is strictly increasing,
while
\begin{equation}
 u(c)=1+\frac2c-m(c)
\end{equation}
is strictly decreasing.  The identities $g(r)=h(0)=0$ and
$g(1/\varphi)=h(1/\varphi)$ give the two endpoints in
\cref{eq:opt:mixed-endpoints}; consequently every
$u\in[\varphi+2,z_*]$ has exactly one parameter pair.

To pass to fixed finite instances, first approximate $(m,t)$ rationally
while retaining the strict cap inequalities, then choose a sufficiently
fine harmonic family, and finally multiply all rational block sizes by a
common integer $H$.  For fixed harmonic depth, quota overshoot, diagonals,
and rounding are $O_u(H)$ whereas the displayed costs are $\Theta(H^2)$.
Send first $H$, then the harmonic depth, and finally the rational
approximation to their limits.  Together with
\cref{lem:opt:mixed-raw-safe}, this proves
\begin{equation}
 V(u)\ge C_{\rm mix}(u)
 \qquad(\varphi+2\le u\le z_*).
 \label{eq:opt:mixed-lower-final}
\end{equation}
The construction above applies directly in the open interval.  At
$u=\varphi+2$ and $u=z_*$ the same statement follows by approaching from
the interior, so the endpoints are included.

\subsection{The large-upper-limit lower bound by cap-free transfer}

\begin{lemma}[Bounded cap-free transfer]
\label{lem:opt:obligatory-transfer}
Let $u\ge1$, and let $V_\infty(P)$ denote the deterministic
size-asymptotic value in the cap-free limit, restricted to
$0\le p_j\le P$.  Then
\begin{equation}
 V(u)\ge V_\infty(u-1).
 \label{eq:opt:obligatory-transfer}
\end{equation}
In particular,
\begin{equation}
 V(u)\ge R_{\rm obl}
 \qquad
 \left(u\ge1+\frac1r=2.752620747896442\ldots\right).
 \label{eq:opt:high-lower}
\end{equation}
\end{lemma}

\begin{proof}
Given a finite-cap algorithm, construct a cap-free algorithm on instances
with $p_j\le u-1$ by running a virtual copy of the finite-cap algorithm.  Copy
ordinary tests and processing actions.  When the virtual algorithm runs an
untouched job raw for $u$ time, physically test and immediately process it in
$1+p_j\le u$ time, hide the answer from the virtual copy, and wait for the
virtual clock.  Every physical completion is no later than its virtual
completion.  Moreover, $\min\{u,1+p_j\}=1+p_j$, so finite-cap and cap-free
optima coincide.  Thus we obtain the following game-value comparison.  For
every finite-cap algorithm
$\mathcal A$, the simulation produces a cap-free algorithm $\mathcal B$ and,
pointwise on every instance with $p_j\le u-1$,
\[
 \ALG_{\mathcal A}\ge\ALG_{\mathcal B},
 \qquad \OPT_u=\OPT_\infty.
\]
Consequently
\[
 \CR_\infty(\mathcal A)
 \ge \CR_\infty(\mathcal B; p\le u-1)
 \ge V_\infty(u-1).
\]
Taking the infimum over $\mathcal A$ proves
\cref{eq:opt:obligatory-transfer}.

The harmonic cap-free family attaining $R_{\rm obl}$ has limiting largest
processing time
\begin{equation}
 1+\log(1+y_*)
 =1+\frac12\log z_*
 =\frac{z_*-1}{2}=\frac1r,
 \qquad y_*=\sqrt{z_*}-1.
\end{equation}
Approximation from below also gives the endpoint, proving
\cref{eq:opt:high-lower}.
\end{proof}

\subsection{Completion of the curve and consistency at the joins}

The lower-bound side of \cref{thm:optional-curve} is now complete:
\cref{prop:opt:binary-lower} supplies the first four branches,
\cref{eq:opt:mixed-lower-final} the fifth, and
\cref{eq:opt:high-lower} the last.  The raw algorithm,
\cref{lem:opt:ute-endpoint-game,lem:opt:zero-prefix}, supplies the matching
upper bounds through $\varphi+2$.  The remaining upper bounds are proved in
\cref{sec:optional-upper}.

For clarity, all five joins are exact:
\begin{align*}
 \rho_I(u_\diamond)&=u_\diamond,\\
 \rho_I(u_0)&=1+\frac1{\sqrt{s_0}}
              =1.682327803828019\ldots,\\
 1+\frac1{\sqrt{(\varphi+2)-1}}&=\varphi
              =C_{\rm mix}(\varphi+2),\\
 C_{\rm mix}(z_*)&=1+r=R_{\rm obl}.
\end{align*}
At $u=1$, both initial branches equal one.  Representative values of the
new mixed branch are
\begin{center}
\begin{tabular}{@{}rrrr@{}}
\toprule
$u$ & $V(u)$ & $m$ & $t$\\
\midrule
$3.800000$ & $1.598851588$ & $0.539725635$ & $1.138984621$\\
$4.000000$ & $1.583283633$ & $0.428863572$ & $1.330517832$\\
$4.071679486$ & $1.579204806$ & $0.381330539$ & $1.412236296$\\
$4.300000$ & $1.571692026$ & $0.198387083$ & $1.734651379$\\
\bottomrule
\end{tabular}
\end{center}

\begin{remark}[Asymptotic scope]
\label{rem:opt:curve-audit}
All statements above concern the fixed-$u$ size-asymptotic ratio.  Below the
plateau, the upper bounds have additive error $O_u(n)$.  On
$u\in[z_*,\infty]$ the same algorithm has one uniform $O(n)$ error.  We do
not claim a finite-instance endpoint inequality with one additive constant.
\end{remark}

The next subsection, \cref{sec:optional-upper}, proves the matching
middle- and high-$u$ upper bounds and thereby completes the proof of
\cref{thm:optional-curve}.
